% --- 1.3.tex ---
\subsection{Trích xuất quan hệ phụ thuộc (Dependency Relation Extraction)}

\subsubsection{Mục tiêu bài toán}
Mục tiêu của tác vụ này là xác định các thành phần ngữ pháp cốt lõi trong một mệnh đề để hiểu được cấu trúc "Ai làm gì" hoặc "Cái gì như thế nào". Thay vì chỉ dừng lại ở việc liệt kê từ loại, hệ thống cần trích xuất được mối quan hệ giữa Động từ chính (Root/Verb) và các thành phần phụ thuộc trực tiếp như Chủ ngữ (Subject), Tân ngữ (Object) và Bổ ngữ (Modifier). Điều này giúp chuẩn hóa dữ liệu cho các tác vụ phân tích ý nghĩa sâu hơn.

\subsubsection{Phương pháp tiếp cận và Triển khai}
Nhóm sử dụng bộ công cụ \texttt{py\_vncorenlp} với tính năng \texttt{parse} (Dependency Parsing) để xây dựng cây cú pháp cho từng mệnh đề. Quy trình triển khai cụ thể như sau:

\begin{enumerate}
    \item \textbf{Xác định Động từ chính (Root Verb):} 
    Hệ thống duyệt qua các từ trong mệnh đề để tìm token có nhãn phụ thuộc là \texttt{root} và từ loại là Động từ (\texttt{V}). Đây được coi là "trung tâm" của quan hệ phụ thuộc.

    \item \textbf{Trích xuất các thành phần phụ thuộc (\texttt{extract\_dependency\_relations}):} 
    Dựa trên chỉ số (\texttt{index}) của Động từ chính, thuật toán tìm kiếm tất cả các từ có quan hệ trực tiếp (nhãn \texttt{head} trùng với \texttt{index} của root). Các nhãn phụ thuộc được ánh xạ sang vai trò ngữ nghĩa tiếng Việt như sau:
    \begin{itemize}
        \item \texttt{sub} $\rightarrow$ \textbf{Chủ ngữ}: Thực thể thực hiện hành động.
        \item \texttt{dob} $\rightarrow$ \textbf{Tân ngữ}: Thực thể chịu tác động của hành động.
        \item \texttt{vmod} $\rightarrow$ \textbf{Bổ ngữ}: Các thông tin bổ sung về thời gian, địa điểm hoặc cách thức.
    \end{itemize}

    \item \textbf{Cấu trúc hóa dữ liệu:} 
    Kết quả được tổ chức dưới dạng một từ điển (dictionary) chứa các thành phần: \textit{Chủ ngữ, Động từ chính, Tân ngữ, Bổ ngữ}. Nếu một thành phần không xuất hiện trong câu, giá trị tương ứng sẽ được để trống.
\end{enumerate}

\subsubsection{Kết quả đầu ra}
Hệ thống xử lý tệp \texttt{clauses.txt} và lưu kết quả vào tệp \texttt{dependency\_relations.txt}. Mỗi dòng kết quả thể hiện cấu trúc ngữ pháp tường minh của một mệnh đề pháp lý. Việc trích xuất thành công các bộ (Subject-Verb-Object) là bước đệm quan trọng để chuyển đổi các điều khoản hợp đồng từ ngôn ngữ tự nhiên sang dạng cấu trúc dữ liệu máy tính có thể hiểu và kiểm tra logic.